\documentclass[12pt]{article}

\usepackage[a4paper, top=2cm, bottom=2cm, left=2.5cm, right=2.5cm]{geometry}

\usepackage{amsmath, amssymb}
\usepackage{tikz}
\usetikzlibrary{arrows.meta,positioning,fit,calc}
\usepackage{multirow}
\usepackage{array}
\usepackage{siunitx}
\usepackage{enumitem}
\usepackage{float}
\usepackage{fontspec}
\usetikzlibrary{decorations.pathmorphing}
\usepackage{pdflscape}
\usepackage{tabularx}
\usepackage{booktabs}
\usepackage{graphicx}
\usepackage{listings}
\usepackage[table]{xcolor}
\usepackage{hyperref}
\usepackage{bookmark}
\usepackage{pdfpages}

% ---------------- FONTS ----------------
\setmainfont[
    UprightFont = Handlee-Regular.ttf,
    AutoFakeBold = 2.5
]{Handlee}

\newfontfamily\headingfont{PT.ttf}
\setmonofont{DejaVu Sans Mono}

% ---------------- HANDWRITTEN MATH ----------------
\usepackage{mathastext}
\Mathastext
\everymath{\displaystyle}

% ---------------- GENERAL ----------------
\setlength{\parindent}{0pt}
\pagenumbering{gobble}
\linespread{1.2}

\newcommand{\mychapter}[2]{%
    \phantomsection
    \hypertarget{#2}{}%
    \bookmark[level=0,dest=#2]{#1}%
    \begin{center}
        {\headingfont\Huge #1}
        \par
        \vspace{4pt}
        \hrule
    \end{center}
}

\newcommand{\mysection}[3]{%
    \phantomsection
    \hypertarget{#3}{}%
    \bookmark[level=1,dest=#3]{#1\space #2}%
    \newpage
    \noindent{\headingfont\Large #1\space #2\par}
    \vspace{4pt}
    \hrule
    \vspace{15pt}
}

\newcommand{\mysubsection}[2]{%
    \vspace{5px}
    \noindent{\headingfont\large #1\space #2\par}
    \vspace{4pt}
    \hrule
    \vspace{15pt}
}

\newcommand{\insertfigure}[2]{%
    \begin{figure}[H]
        \centering
        \fbox{\includegraphics[width=#1\textwidth]{../2. Figures/#2}}
    \end{figure}
}

\begin{document}

\mychapter{Chapter 6 : Second Law of Thermodynamics}{6f4y5kdh}

\begin{center}
\renewcommand{\arraystretch}{1.5}
\setlength{\tabcolsep}{10pt}

\begin{tabular}{|c|p{0.78\textwidth}|}
\hline
\textbf{Section} & \textbf{Core Revision Focus} \\
\hline

6.1 &
Introduces the second law as the basis for determining process direction, feasibility, and the limitations of energy conversion. \\
\hline

6.2 &
Defines thermal energy reservoirs and distinguishes heat sources from heat sinks used in thermodynamic cycles. \\
\hline

6.3 &
Develops heat-engine operation, thermal efficiency, waste-heat rejection, and the Kelvin--Planck statement. \\
\hline

6.4 &
Covers refrigerators and heat pumps, COP, energy-efficiency measures, the Clausius statement, and their equivalence. \\
\hline

6.5 &
Explains perpetual-motion machines of the first and second kinds and the laws of thermodynamics they violate. \\
\hline

6.6 &
Distinguishes reversible and irreversible processes and identifies internal, external, total reversibility, and irreversibilities. \\
\hline

6.7 &
Develops the reversible Carnot cycle, its four processes, heat/work interactions, and maximum possible efficiency. \\
\hline

6.8 &
Establishes the Carnot principles: reversible engines are most efficient and all reversible engines between the same reservoirs have equal efficiency. \\
\hline

6.9 &
Develops the thermodynamic temperature scale and the relation between heat-transfer ratios and absolute temperatures. \\
\hline

6.10 &
Applies Carnot efficiency to heat engines and examines temperature limits, efficiency improvement, and energy quality/degradation. \\
\hline

6.11 &
Develops Carnot refrigerator and heat-pump COP relations and examines the effect of temperature limits on performance. \\
\hline

\end{tabular}
\end{center}

% ============================================================
% 6.1
% ============================================================

\mysection{6.1}{Introduction to the Second Law}{6.1}

\mysubsection{}{Core Idea}

The first law gives \textbf{energy conservation}; the second law gives the \textbf{direction and limitations of processes}.

\begin{equation*}
\text{Feasible Process}
\Longrightarrow
\text{First Law Satisfied}
+
\text{Second Law Satisfied}
\end{equation*}

\mysubsection{}{Limitations of the First Law}

\begin{itemize}[leftmargin=*]
    \item First law determines the \textbf{quantity} of energy and energy transformations.
    \item It does not determine the \textbf{direction} of a process.
    \item Natural processes have preferred directions; spontaneous reverse processes generally do not occur.
    \item The second law provides the criterion for process direction and theoretical performance limits.
\end{itemize}

\mysubsection{}{Energy Quality}

\begin{itemize}[leftmargin=*]
    \item First law $\rightarrow$ energy quantity/conservation.
    \item Second law $\rightarrow$ energy quality, work potential, degradation and process limitations.
    \item For equal quantities of energy, higher-temperature energy generally has greater work-producing potential.
\end{itemize}

\begin{equation*}
\text{Second Law}
\longrightarrow
\text{Direction}
+
\text{Energy Quality}
+
\text{Performance Limits}
\end{equation*}

% ============================================================
% 6.2
% ============================================================

\mysection{6.2}{Thermal Energy Reservoirs}{6.2}

\mysubsection{}{Thermal Energy Reservoir}

A \textbf{thermal energy reservoir} is a body that can supply or absorb a finite amount of heat while maintaining essentially constant temperature.

\begin{equation*}
\text{Thermal Energy Capacity}=mc
\end{equation*}

where $m$ = mass and $c$ = specific heat.

\textbf{Reservoir condition:}
the thermal energy capacity is sufficiently large that the heat transfer causes negligible temperature change.

Typical models: oceans, lakes, rivers, atmospheric air, two-phase systems, controlled furnaces.

\mysubsection{}{Source and Sink}

\begin{equation*}
\text{Source}\longrightarrow\text{Heat}
\end{equation*}

A \textbf{source} supplies heat.

\begin{equation*}
\text{Heat}\longrightarrow\text{Sink}
\end{equation*}

A \textbf{sink} absorbs heat.

% ============================================================
% 6.3
% ============================================================

\mysection{6.3}{Heat Engines}{6.3}

\mysubsection{}{Heat Engine}

A \textbf{heat engine} is a cyclic device that receives heat from a high-temperature source, converts part of it into work, and rejects the remainder to a low-temperature sink.

\begin{equation*}
Q_H=W_{\mathrm{net,out}}+Q_L
\end{equation*}

where

\begin{equation*}
Q_H=\text{heat supplied},\qquad
Q_L=\text{heat rejected}
\end{equation*}

\begin{equation*}
T_H>T_L
\end{equation*}

\mysubsection{}{Thermal Efficiency}

\begin{equation*}
\eta_{\mathrm{th}}
=
\frac{W_{\mathrm{net,out}}}{Q_H}
\end{equation*}

Using the cycle energy balance,

\begin{equation*}
W_{\mathrm{net,out}}=Q_H-Q_L
\end{equation*}

hence

\begin{equation*}
\boxed{
\eta_{\mathrm{th}}
=
1-\frac{Q_L}{Q_H}
}
\end{equation*}

\begin{equation*}
0\leq\eta_{\mathrm{th}}<1
\end{equation*}

\mysubsection{}{Kelvin--Planck Statement}

\textbf{A cyclic device cannot receive heat from a single reservoir and produce a net amount of work.}

Therefore, a cyclic heat engine must interact with both a high-temperature and a low-temperature reservoir.

\begin{equation*}
\boxed{\eta_{\mathrm{th}}<1}
\end{equation*}

Complete conversion of heat from a single reservoir into cyclic work is impossible.

\mysubsection{}{Why Heat Rejection Is Necessary}

For continuous cyclic operation,

\begin{equation*}
Q_H=W_{\mathrm{net,out}}+Q_L
\end{equation*}

Thus some heat must be rejected to the low-temperature reservoir. Waste heat rejection is a \textbf{fundamental second-law requirement}, not merely a consequence of friction or imperfect construction.

% ============================================================
% 6.4
% ============================================================

\mysection{6.4}{Refrigerators and Heat Pumps}{6.4}

\mysubsection{}{Refrigerator}

A refrigerator transfers heat from a low-temperature region to a high-temperature region using work input.

\begin{equation*}
W_{\mathrm{net,in}}=Q_H-Q_L
\end{equation*}

where $Q_L$ is the refrigeration effect.

\begin{equation*}
\boxed{
\mathrm{COP}_R
=
\frac{Q_L}{W_{\mathrm{net,in}}}
=
\frac{Q_L}{Q_H-Q_L}
=
\frac{1}{Q_H/Q_L-1}
}
\end{equation*}

\textbf{Note:} $\mathrm{COP}_R$ may exceed unity; COP is not thermal efficiency.

\mysubsection{}{Heat Pump}

A heat pump uses work to transfer heat from a low-temperature source to a high-temperature space.

\begin{equation*}
\boxed{
\mathrm{COP}_{HP}
=
\frac{Q_H}{W_{\mathrm{net,in}}}
=
\frac{Q_H}{Q_H-Q_L}
=
\frac{1}{1-Q_L/Q_H}
}
\end{equation*}

\mysubsection{}{Refrigerator vs Heat Pump}

\begin{center}
\begin{tabular}{|c|c|c|}
\hline
\textbf{Feature} & \textbf{Refrigerator} & \textbf{Heat Pump} \\
\hline
Purpose & Cooling & Heating \\
\hline
Desired effect & $Q_L$ & $Q_H$ \\
\hline
Input & $W_{\mathrm{net,in}}$ & $W_{\mathrm{net,in}}$ \\
\hline
Performance & $\mathrm{COP}_R$ & $\mathrm{COP}_{HP}$ \\
\hline
\end{tabular}
\end{center}

Since

\begin{equation*}
Q_H=Q_L+W_{\mathrm{net,in}}
\end{equation*}

\begin{equation*}
\boxed{
\mathrm{COP}_{HP}
=
\mathrm{COP}_R+1
}
\end{equation*}

\mysubsection{}{EER and SEER}

\begin{equation*}
\mathrm{EER}
=
\frac{\text{rate of heat removal}}
{\text{rate of electricity consumption}}
\qquad [\mathrm{Btu/Wh}]
\end{equation*}

\begin{equation*}
\mathrm{SEER}
=
\frac{\text{seasonal heat removed}}
{\text{seasonal electricity consumed}}
\qquad [\mathrm{Btu/Wh}]
\end{equation*}

\begin{equation*}
\boxed{
\mathrm{EER}=3.412\,\mathrm{COP}_R
}
\end{equation*}

\mysubsection{}{Clausius Statement}

\textbf{A cyclic device cannot produce no effect other than transferring heat from a lower-temperature body to a higher-temperature body.}

A refrigerator can transfer heat from cold to hot only with external work input.

\mysubsection{}{Kelvin--Planck vs Clausius}

The two statements are \textbf{equivalent consequences} of the second law.

\begin{equation*}
\text{Violation of Kelvin--Planck}
\Longrightarrow
\text{Violation of Clausius}
\end{equation*}

and conversely.

% ============================================================
% 6.5
% ============================================================

\mysection{6.5}{Perpetual-Motion Machines}{6.5}

\mysubsection{}{Types}

\begin{center}
\begin{tabular}{|c|c|c|}
\hline
\textbf{Machine} & \textbf{Law Violated} & \textbf{Violation} \\
\hline
PMM1 & First law & Creates energy \\
\hline
PMM2 & Second law & Violates process direction/energy-conversion limits \\
\hline
\end{tabular}
\end{center}

\mysubsection{}{PMM1}

\textbf{Perpetual-motion machine of the first kind} violates conservation of energy.

\begin{equation*}
\text{PMM1}\Longrightarrow\text{Energy Creation}
\end{equation*}

\mysubsection{}{PMM2}

A \textbf{PMM2} violates the second law without necessarily violating energy conservation.

A cyclic device that:

\begin{itemize}[leftmargin=*]
    \item receives heat from a single reservoir, and
    \item produces net work
\end{itemize}

is a PMM2.

\mysubsection{}{Essential Distinction}

\begin{equation*}
\text{First Law}\rightarrow\text{Conservation}
\end{equation*}

\begin{equation*}
\text{Second Law}\rightarrow\text{Direction and Limitations}
\end{equation*}

% ============================================================
% 6.6
% ============================================================

\mysection{6.6}{Reversible and Irreversible Processes}{6.6}

\mysubsection{}{Reversible Process}

A process is \textbf{reversible} if the system and surroundings can both be restored completely to their initial states with no net effect.

For the combined forward and reverse processes,

\begin{equation*}
Q_{\mathrm{net}}=0,
\qquad
W_{\mathrm{net}}=0
\end{equation*}

\textbf{System restoration alone is insufficient; the surroundings must also be restored.}

\mysubsection{}{Irreversible Process}

An \textbf{irreversible process} cannot be reversed without leaving a net change in the system or surroundings.

\textbf{All natural processes are irreversible.}

Reversible processes are ideal limiting models.

\mysubsection{}{Engineering Importance}

For work-producing devices:

\begin{equation*}
\boxed{
\text{Reversible}\rightarrow
W_{\mathrm{out,max}}
}
\end{equation*}

For work-consuming devices:

\begin{equation*}
\boxed{
\text{Reversible}\rightarrow
W_{\mathrm{in,min}}
}
\end{equation*}

Reversible operation therefore establishes the theoretical performance limit.

\mysubsection{}{Major Irreversibilities}

\begin{itemize}[leftmargin=*]
    \item friction,
    \item unrestrained expansion,
    \item mixing,
    \item heat transfer through a finite temperature difference,
    \item electrical resistance,
    \item inelastic deformation,
    \item chemical reactions.
\end{itemize}

\mysubsection{}{Types of Reversibility}

\begin{center}
\begin{tabular}{|c|c|}
\hline
\textbf{Type} & \textbf{Condition} \\
\hline
Internally reversible & No irreversibilities inside system boundary \\
\hline
Externally reversible & No irreversibilities outside system boundary \\
\hline
Totally reversible & No internal or external irreversibilities \\
\hline
\end{tabular}
\end{center}

\textbf{Internally reversible:}
system passes through equilibrium states and forward/reverse paths coincide.

\textbf{Externally reversible:}
heat transfer occurs without a finite temperature difference at the boundary.

\begin{equation*}
\boxed{
\text{Totally reversible}
=
\text{Internally reversible}
+
\text{Externally reversible}
}
\end{equation*}

% ============================================================
% 6.7
% ============================================================

\mysection{6.7}{The Carnot Cycle}{6.7}

\mysubsection{}{Carnot Cycle}

The Carnot cycle is a completely reversible cycle consisting of:

\begin{itemize}[leftmargin=*]
    \item two reversible isothermal processes,
    \item two reversible adiabatic processes.
\end{itemize}

It operates between

\begin{equation*}
T_H>T_L
\end{equation*}

\mysubsection{}{Four Processes}

\begin{center}
\begin{tabular}{|c|c|c|c|}
\hline
\textbf{Process} & \textbf{Type} & \textbf{Heat} & \textbf{Temperature} \\
\hline
$1\rightarrow2$ &
Isothermal expansion &
$Q_H$ in &
$T_H=\mathrm{constant}$ \\
\hline
$2\rightarrow3$ &
Adiabatic expansion &
$Q=0$ &
$T_H\rightarrow T_L$ \\
\hline
$3\rightarrow4$ &
Isothermal compression &
$Q_L$ out &
$T_L=\mathrm{constant}$ \\
\hline
$4\rightarrow1$ &
Adiabatic compression &
$Q=0$ &
$T_L\rightarrow T_H$ \\
\hline
\end{tabular}
\end{center}

\mysubsection{}{Cycle Relations}

\begin{equation*}
1\rightarrow2\rightarrow3\rightarrow4\rightarrow1
\end{equation*}

\begin{equation*}
\boxed{
W_{\mathrm{net,out}}=Q_H-Q_L
}
\end{equation*}

On the $P$--$V$ diagram, the enclosed area represents net work output:

\begin{equation*}
W_{\mathrm{net,out}}
=
W_{\mathrm{expansion}}
-
W_{\mathrm{compression}}
\end{equation*}

\mysubsection{}{Maximum Performance}

Because the Carnot cycle is completely reversible,

\begin{equation*}
\boxed{
\eta_{\mathrm{Carnot}}
=
\eta_{\mathrm{rev}}
=
\eta_{\mathrm{max}}
}
\end{equation*}

for a heat engine operating between the same two temperature limits.

No actual engine operating between the same limits can exceed Carnot efficiency.

\mysubsection{}{Reversed Carnot Cycle}

Reversing every process gives a Carnot refrigerator/heat pump:

\begin{equation*}
Q_L:\text{absorbed},\qquad
Q_H:\text{rejected},\qquad
W_{\mathrm{net,in}}:\text{supplied}
\end{equation*}

% ============================================================
% 6.8
% ============================================================

\mysection{6.8}{The Carnot Principles}{6.8}

\mysubsection{}{First Carnot Principle}

\textbf{An irreversible heat engine is less efficient than a reversible heat engine operating between the same two reservoirs.}

\begin{equation*}
\boxed{
\eta_{\mathrm{th,irrev}}
<
\eta_{\mathrm{th,rev}}
}
\end{equation*}

\begin{equation*}
\text{Irreversibilities}
\longrightarrow
\text{Lower efficiency}
\end{equation*}

\mysubsection{}{Second Carnot Principle}

\textbf{All reversible heat engines operating between the same two reservoirs have the same efficiency.}

\begin{equation*}
\boxed{
\eta_{\mathrm{rev,1}}
=
\eta_{\mathrm{rev,2}}
=
\eta_{\mathrm{rev,3}}
}
\end{equation*}

Their efficiency depends only on the two reservoir temperatures, not on:

\begin{itemize}[leftmargin=*]
    \item working fluid,
    \item engine design,
    \item method of executing the cycle.
\end{itemize}

\mysubsection{}{Combined Result}

For engines operating between the same two reservoirs,

\begin{equation*}
\boxed{
\eta_{\mathrm{irrev}}
<
\eta_{\mathrm{rev,1}}
=
\eta_{\mathrm{rev,2}}
=
\eta_{\mathrm{rev,3}}
}
\end{equation*}

The reversible engine establishes the \textbf{upper performance limit}.

% ============================================================
% 6.9
% ============================================================

\mysection{6.9}{The Thermodynamic Temperature Scale}{6.9}

\mysubsection{}{Basic Result}

For reversible heat engines, efficiency depends only on reservoir temperatures:

\begin{equation*}
\eta_{\mathrm{th,rev}}=g(T_H,T_L)
\end{equation*}

Therefore,

\begin{equation*}
\left(\frac{Q_H}{Q_L}\right)_{\mathrm{rev}}
=
f(T_H,T_L)
\end{equation*}

\mysubsection{}{Three-Engine Relation}

For three reversible engines,

\begin{equation*}
\frac{Q_1}{Q_2}=f(T_1,T_2),
\qquad
\frac{Q_2}{Q_3}=f(T_2,T_3),
\qquad
\frac{Q_1}{Q_3}=f(T_1,T_3)
\end{equation*}

Since

\begin{equation*}
\frac{Q_1}{Q_3}
=
\frac{Q_1}{Q_2}
\frac{Q_2}{Q_3}
\end{equation*}

the functional form can be written as

\begin{equation*}
\boxed{
f(T_1,T_2)
=
\frac{\phi(T_1)}{\phi(T_2)}
}
\end{equation*}

Hence,

\begin{equation*}
\boxed{
\frac{Q_H}{Q_L}
=
\frac{\phi(T_H)}{\phi(T_L)}
}
\]

for a reversible heat engine.

\mysubsection{}{Kelvin Temperature Scale}

Choose

\begin{equation*}
\phi(T)=T
\end{equation*}

Then

\begin{equation*}
\boxed{
\left(\frac{Q_H}{Q_L}\right)_{\mathrm{rev}}
=
\frac{T_H}{T_L}
}
\end{equation*}

or

\begin{equation*}
\boxed{
\frac{Q_H}{Q_L}
=
\frac{T_H}{T_L}
}
\qquad\text{(reversible)}
\end{equation*}

This defines the thermodynamic/Kelvin absolute-temperature scale.

\mysubsection{}{Important Kelvin-Scale Facts}

\begin{itemize}[leftmargin=*]
    \item Absolute temperature scale.
    \item Independent of properties of a particular thermometric substance.
    \item Temperature ratios correspond directly to reversible heat-transfer ratios.
    \item Lower limit is absolute zero.
\end{itemize}

\mysubsection{}{Triple Point and Celsius Scale}

A fixed temperature reference is required to establish the temperature-unit magnitude.

The source notes use the \textbf{triple point of water} as the fixed reference.

\begin{equation*}
1\,\mathrm{K}\equiv1\,^{\circ}\mathrm{C}
\end{equation*}

\begin{equation*}
\boxed{
T(^{\circ}\mathrm{C})
=
T(\mathrm{K})-273.15
}
\end{equation*}

Temperature differences have equal numerical magnitudes in K and $^{\circ}$C.

\mysubsection{}{Absolute Zero}

\begin{equation*}
\boxed{T\geq0}
\end{equation*}

Absolute zero is the lower limit of thermodynamic temperature.

\mysubsection{}{Practical Determination}

Absolute temperature need not be measured with an actual reversible heat engine. A \textbf{constant-volume ideal-gas thermometer} can be used for practical temperature measurement.

% ============================================================
% 6.10
% ============================================================

\mysection{6.10}{The Carnot Heat Engine}{6.10}

\mysubsection{}{Carnot Efficiency}

For a reversible heat engine,

\begin{equation*}
\boxed{
\eta_{\mathrm{th,Carnot}}
=
1-\frac{T_L}{T_H}
}
\end{equation*}

\textbf{Condition:} $T_H$ and $T_L$ must be \textbf{absolute temperatures}.

\begin{equation*}
\eta_{\mathrm{th,actual}}
<
\eta_{\mathrm{th,Carnot}}
\end{equation*}

and, in general,

\begin{equation*}
\boxed{
\eta_{\mathrm{th}}\leq\eta_{\mathrm{th,rev}}
}
\end{equation*}

for a physically possible engine operating between the same temperature limits.

\mysubsection{}{Temperature Effects}

From

\begin{equation*}
\eta_{\mathrm{th,rev}}
=
1-\frac{T_L}{T_H}
\end{equation*}

\begin{equation*}
T_H\uparrow
\quad\Longrightarrow\quad
\eta_{\mathrm{th,rev}}\uparrow
\end{equation*}

\begin{equation*}
T_L\downarrow
\quad\Longrightarrow\quad
\eta_{\mathrm{th,rev}}\uparrow
\end{equation*}

Thus, theoretically:

\begin{equation*}
\boxed{
\text{Increase }T_H
\quad\text{and/or}\quad
\text{decrease }T_L
}
\end{equation*}

to increase the maximum efficiency.

Practical limits arise from material temperature limits and available cooling temperatures.

\mysubsection{}{Dependence of Carnot Efficiency}

Carnot efficiency depends only on

\begin{equation*}
T_H,\quad T_L
\end{equation*}

and not on:

\begin{itemize}[leftmargin=*]
    \item working fluid,
    \item engine design,
    \item details of the reversible cycle,
    \item method of executing reversible processes.
\end{itemize}

\mysubsection{}{Quality of Energy}

The first law concerns \textbf{quantity}; the second law also concerns \textbf{quality/work potential}.

\begin{equation*}
\text{Higher temperature}
\Longrightarrow
\text{Higher work-producing potential}
\end{equation*}

All work can be completely converted to heat, but only part of heat can be converted to work.

\begin{equation*}
\text{Work}\rightarrow\text{Heat}
\quad\text{(complete conversion possible)}
\end{equation*}

\begin{equation*}
\text{Heat}\rightarrow\text{Work}
\quad\text{(second-law limited)}
\end{equation*}

\mysubsection{}{Energy Degradation}

When heat flows from high temperature to low temperature:

\begin{equation*}
\text{Energy quantity remains conserved}
\end{equation*}

but

\begin{equation*}
\text{Work-producing potential decreases}
\end{equation*}

Thus,

\begin{equation*}
\boxed{
\text{Energy degradation}
=
\text{reduction in work-producing potential}
}
\end{equation*}

% ============================================================
% 6.11
% ============================================================

\mysection{6.11}{The Carnot Refrigerator and Heat Pump}{6.11}

\mysubsection{}{Carnot Refrigerator}

For a refrigerator operating reversibly between $T_H$ and $T_L$,

\begin{equation*}
W_{\mathrm{net,in}}=Q_H-Q_L
\end{equation*}

\begin{equation*}
\boxed{
\mathrm{COP}_{R,\mathrm{rev}}
=
\frac{Q_L}{Q_H-Q_L}
}
\end{equation*}

Using

\begin{equation*}
\frac{Q_H}{Q_L}
=
\frac{T_H}{T_L}
\end{equation*}

gives

\begin{equation*}
\boxed{
\mathrm{COP}_{R,\mathrm{rev}}
=
\frac{1}{T_H/T_L-1}
=
\frac{T_L}{T_H-T_L}
}
\end{equation*}

\mysubsection{}{Carnot Heat Pump}

For a reversible heat pump,

\begin{equation*}
\boxed{
\mathrm{COP}_{HP,\mathrm{rev}}
=
\frac{Q_H}{Q_H-Q_L}
}
\end{equation*}

and therefore

\begin{equation*}
\boxed{
\mathrm{COP}_{HP,\mathrm{rev}}
=
\frac{1}{1-T_L/T_H}
=
\frac{T_H}{T_H-T_L}
}
\end{equation*}

\mysubsection{}{Temperature Effects on Refrigerator COP}

\begin{equation*}
\mathrm{COP}_{R,\mathrm{rev}}
=
\frac{T_L}{T_H-T_L}
\end{equation*}

At fixed $T_H$,

\begin{equation*}
T_L\downarrow
\quad\Longrightarrow\quad
\mathrm{COP}_{R,\mathrm{rev}}\downarrow
\end{equation*}

and as

\begin{equation*}
T_L\rightarrow0,
\qquad
\mathrm{COP}_R\rightarrow0
\end{equation*}

Hence, refrigeration becomes increasingly difficult as the refrigerated space approaches absolute zero.

\mysubsection{}{Temperature Effects on Heat-Pump COP}

\begin{equation*}
\mathrm{COP}_{HP,\mathrm{rev}}
=
\frac{T_H}{T_H-T_L}
\end{equation*}

A larger temperature lift,

\begin{equation*}
T_H-T_L\uparrow
\end{equation*}

requires more work and therefore

\begin{equation*}
\mathrm{COP}_{HP,\mathrm{rev}}\downarrow
\end{equation*}

Hence, heat-pump performance improves when the required temperature difference is reduced.

\mysubsection{}{Relationship Between COPs}

\begin{equation*}
\boxed{
\mathrm{COP}_{HP}
=
\mathrm{COP}_R+1
}
\end{equation*}

and for reversible operation,

\begin{equation*}
\boxed{
\mathrm{COP}_{HP,\mathrm{rev}}
=
\mathrm{COP}_{R,\mathrm{rev}}+1
}
\end{equation*}

Therefore,

\begin{equation*}
\mathrm{COP}_{HP}>\mathrm{COP}_R
\end{equation*}

for the same operating conditions.

\mysubsection{}{Reversible vs Irreversible Devices}

For the same temperature limits:

\begin{center}
\begin{tabular}{|c|c|}
\hline
\textbf{Device} & \textbf{Second-Law Limit} \\
\hline
Irreversible refrigerator &
$\mathrm{COP}_R<\mathrm{COP}_{R,\mathrm{rev}}$ \\
\hline
Reversible refrigerator &
$\mathrm{COP}_R=\mathrm{COP}_{R,\mathrm{rev}}$ \\
\hline
Impossible refrigerator &
$\mathrm{COP}_R>\mathrm{COP}_{R,\mathrm{rev}}$ \\
\hline
Irreversible heat pump &
$\mathrm{COP}_{HP}<\mathrm{COP}_{HP,\mathrm{rev}}$ \\
\hline
Reversible heat pump &
$\mathrm{COP}_{HP}=\mathrm{COP}_{HP,\mathrm{rev}}$ \\
\hline
Impossible heat pump &
$\mathrm{COP}_{HP}>\mathrm{COP}_{HP,\mathrm{rev}}$ \\
\hline
\end{tabular}
\end{center}

The reversible values are the \textbf{maximum possible COPs} between the specified temperature limits.

\end{document}